\scanpage{225}
% archive sheet number: 224

% slanted
\margin{\deleted{$X$ \inserted{\unsure{normal irréduct.}} de type fini sur un \unsure{corps}}}

On a vu plus haut qu'une place est très loin d'être caractérisée par le foncteur fibre correspondant, \deleted{Mais} considéré comme foncteur dans les $\mathbb{Z}$-modules. Mais que se passe-t-il si on le regarde comme foncteur dans les structures de Hodge ? \deleted{Si on part} \inserted{Cela \deleted{\unsure{doit être}} est vrai pour un schéma \unsure{st local|strict local}, et pour la \unsure{même raison}}

% slanted
\margin{attention, il faut parler des \unsure{pts complexes|pl. complexes}, \unsure{ou|et non} seulement des \unsure{pl. ré. conj.}, pour pouvoir définir la structure de Hodge \unsure{---|...}}

\diagram{\begin{tikzcd}[ampersand replacement=\&,column sep=small,row sep=small,execute at end picture={
\draw[thin] (topleft.north east) -- (bottomleft.south west);
\draw[thin] ($(topleft.east)+(.07,0)$) -- ($(bottomleft.east)+(-.02,-.08)$);
\draw[thin] (bottommid.north) .. controls +(-.5,-.8) and +(-.4,.9) .. (bottomright.west);
}]
\& |[alias=topleft]|\deleted{\unsure{X_i|X_0}}\arrow[r,dash]\arrow[d,dash]\arrow[dl,dash] \& X\arrow[d,dash]\arrow[dl,dash]\\
|[alias=bottomleft]|\deleted{\unsure{\mathbb{Q}|k_0}}\arrow[r,dash] \& |[alias=bottommid]|\deleted{\unsure{A_i}}\arrow[r,dash] \& |[alias=bottomright]|k
\end{tikzcd}}

\diagram{\begin{tikzcd}[ampersand replacement=\&,row sep=small] X\arrow[d,dash]\\ {}\end{tikzcd}}

\deleted{\unsure{par-dessus le point}, on est ramené au cas où \unsure{deux places coïncident déjà sur} $k_0=$ clôture alg. de $\unsure{\mathbb{Q}|Q}$ dans $k$, et puis au cas où $X$ est de type}

\deleted{$X$ est de type fini sur \deleted{\unsure{un}} $\mathbb{Q}$, puis par \unsure{changem\textsuperscript{t}} de base au cas où $X$ est connexe de type fini sur $\mathbb{C}$, et qu'on a deux pts \inserted{ordinaires} de $X$ à val. dans $\mathbb{C}$.}

\deleted{Supposant pour simplifier que \deleted{\unsure{la courbe}} $X$ est \deleted{\unsure{réunion des}\ \illegible{1}} \deleted{affine,} quasi-projectif, \deleted{\unsure{est affine}}, on est ramené au cas où on peut supposer $X=P^{r}_{\mathbb{C}}$.}

Il semble qu'on se ramène facilement au cas de $X=\mathbb{P}^{\unsure{*|1}}_{\mathbb{C}}-(0)-(1)-(\infty)$, \inserted{et deux pts à val. ds $\mathbb{C}$ ord.,} où la théorie \unsure{classique} \deleted{\unsure{doit}} le donner --- [prendre comme motif celui défini par un schéma en courbes elliptiques sur $X$, d'invariant l'indéterminée $j$].
